\documentclass[submitted,11pt]{article}
\usepackage[utf8]{inputenc}
\usepackage{els}
\usepackage{ii}
\usepackage{rebuttal}
\newcommand{\RT}{Reviewer~\#2\xspace} 
\newcommand{\RO}{Reviewer~\#1\xspace}

\newcommand{\TDE}{Transverse Doppler Effect\xspace}
\newcommand{\tDe}{transverse Doppler effect\xspace}
\newcommand{\EP}{\Ehr paradox\xspace}
\newcommand{\mspp}[1]{\mkbibbrackets{p.~#1}}
\newcommand{\igamma}{\ensuremath{\gamma^{-1}}\xspace}

\title{Response to the Referees for SYNT-D-26-00214}

\begin{document}
\maketitle
\section*{Response to reviewers' comments}

\noindent I would like to thank the reviewers for their time and effort in reviewing the manuscript. The paper has been revised carefully in light of their comments.

\noindent The principal revisions concern three issues:
\begin{enumerate}
\item the distinction between the \s{real} and \s{apparent} character of relativistic effects, following Einstein's terminology;
\item the clarification of the sense in which the transverse Doppler effect \s{vanishes} for a co-moving observer;
\item the role and philosophical significance of Bell-style dynamical explanations in relation to Einstein's kinematical approach.
\end{enumerate}


\noindent All changes have been tracked and are color-coded for your convenience. A version of the paper without track changes has been included for comparison.

%3. If you do those, your response letter can honestly say the paper has undergone substantive revision rather than only clarification.

\reviewersection

\reviewersection

\begin{point}
First, it is claimed that length contraction and time dilation, in the eyes of Einstein, are "apparent" effects because they "vanish" for the co-moving observer, but are nonetheless "real" since they "cannot be eliminated in all frames at once" (pp. 2, 22, footnote 8). This latter criterion of reality is almost surely NOT what Einstein had in mind. 
\end{point}

%The nontrivial premise is not the invariance of \(N\) as such, but the invariance of the proper frequency \(N/\Delta t'\).

\begin{reply}
My explanation of  the distinction between \emph{real} and \emph{apparent} was unclear. It is introduced as an exegetical reconstruction of Einstein's own usage\todo{quotes}. The clearest formulation occurs in his 1915 letter to Lorentz:

\begin{quote}
Regarding the erroneous view that the Lorentz contraction was \s{merely apparent,} \origins{scheinbar} I am not free from guilt, without ever having myself lapsed into that error. It is real \origins{wirklich}, i.e., measurable with \rac, and at the same time apparent \origins{scheinbar} to the extent that it is not present for the co-moving observers \lettercpaep{Einstein}{Lorentz}{23}{1}{1915}[8][47].
\end{quote}
%
Accordingly, the paper does not claim that an effect is real because it \emph{cannot be eliminated in all frames at once} (this applies to cases like the \Ehr paradox, see \ref{global}). Rather, it reconstructs Einstein's own characterization: 

%an effect is \emph{real} insofar as it is experimentally measurable, yet \emph{apparent} insofar as it disappears for co-moving observers.

\begin{itemize}
\item apparent = absent for the co-moving observer.
\item real = experimentally measurable,
\end{itemize}
%
The transverse Doppler effect illustrates precisely this point. A spectrometer co-moving with the emitting ion records its proper frequency, so that no shift is observed in that experimental arrangement. A spectrometer in relative motion, by contrast, measures the shifted frequency. The effect is therefore relational, but this is precisely what Einstein meant by calling it simultaneously \emph{real} and \emph{apparent}. The fact that the shift disappears in one observational arrangement does not make it illusory; it remains experimentally measurable in another (see also my response to \ref{vanish}).

As I mentioned in a footnote, the distinction may also be illustrated by the familiar case of stellar aberration. The displaced position of the star is \emph{apparent}, since it depends on the observer's state of motion. Yet aberration is \emph{real} in an operational sense. If an astronomer points a telescope toward the star's true position rather than its aberrated position, the starlight will strike the wall of the telescope tube instead of reaching the eyepiece.
\end{reply}



\begin{point}
Applications of length contraction and time dilation can, in the appropriate circumstances, involve non-classical empirical predictions, such as the breaking of the string in Bell's two-rockets thought-experiment (footnote 1),  stresses in a rotating disk (p. 2), the experimental confirmation of the transverse Doppler shift (section 2), and time retardation in the twin's "paradox" scenario (not mentioned in the paper). These predictions of special relativity are ones that ALL observers will agree on, including the observer in the relevant co-moving frame(s). It is in this sense that length contraction and time dilation are
supposedly real. Indeed in footnote 1, the author explicitly states that the examples of "stresses and the possible structural failure of the rope (or of the rotating disk)" mean that length contraction is physically real for Einstein — these are coordinate-independent predictions. \label{global}
\end{point}

\begin{reply}
I agree that the original formulation did not sufficiently distinguish the local and global senses in which a relativistic effect may be said to \s{disappear} for a co-moving observer.

In the local case, a particular length contraction or time dilation can be made to disappear by passing to the rest frame of the relevant material element or clock. In that operational sense, the effect is \s{apparent}: the co-moving observer measures the proper length, proper time, or proper frequency. It is nevertheless \s{real} in Einstein's sense because it is measurable using rods and clocks in a non-co-moving frame. The \tDe should be understood in this context.

In extended or non-inertial situations, however, the local rest frames need not fit together into a single global co-moving inertial frame. This is what happens in the rotating disk: each element of the rim has an instantaneous rest frame, but these frames cannot be combined into a single global rest frame for the disk. In the original paper, I mentioned the Ehrenfest disk not as a general criterion for the reality of Lorentz contraction, but to recall the historical context in which the debate over whether Lorentz contraction is \s{real} or merely \s{apparent} first became explicit. In this correspondence, Einstein invokes the rotating disk to show that Lorentz contraction is a kinematical effect and therefore \s{apparent} in the sense that it is absent for suitably co-moving observers. At the same time, it is \s{real} in the sense that such contractions cannot be made to disappear globally for all inertial observers simultaneously, a fact manifested by the stresses that arise in the rotating disk \lettercpaep{Einstein}{\Var}{3}{3}{1911}[5{[}10{]}][257a].

I have now removed the reference to the \Ehr paradox, as I agree that it was potentially misleading.


%The transverse Doppler effect occupies an important place in the paper because, unlike the early discussions of length contraction, it offered Einstein a concrete empirical route for testing time dilation. A co-moving spectrograph records the proper frequency of the emitting ion, so the shift disappears in that local observational arrangement. A non-co-moving spectrograph, however, records the displaced spectral line. This made time dilation experimentally accessible in a way that length contraction, at the time, was not. 


%I agree with the referee remark. I should have distinguished more clearly between two cases: \begin{inparaenum}[(a)] \item In the standard inertial case, the contraction or dilation disappears (see \ref{vanish}) in the corresponding co-moving frame. \item In more complex situations, such as the rotating disk or the twin paradox, no single global co-moving inertial frame exists.
%\end{inparaenum}
% 
%% his is precisely why Einstein regarded them as \s{real}. However, this does not conflict with his simultaneous characterization of them as \s{apparent}. 
% 
%These examples do not undermine the distinction between real and apparent; rather, they illustrate why Einstein regarded the two characterizations as compatible. Historically,  to illustrate precisely the interplay between the apparent and the real character of Lorentz contraction Einstein used case (b) that is the \Ehr paradox. The reason is probably that there is no direct testable of empirical consequence of Lorentz contraction alone.  
%
%
%%In uniform translation, the description in the co-moving frame and the description in a non-co-moving inertial frame remain globally equivalent: all collinear parts of the rod acquire the same Lorentz factor, so that the ratio between the whole and its parts is preserved. In rotation this equivalence breaks down.
%
%There is however a direct empirical consequence of time dilation, namely the \TD. Possibly for this reason, Einstein the twin paradox is not employed by Einstein in this context. Nevertheless, from a conceptual standpoint it can be viewed as the temporal analogue of the Ehrenfest argument: just as no global co-moving inertial frame exists for the rotating disk, none exists for the accelerated twin. In both cases, the local disappearance of the relativistic effect cannot be globalized, and observable, frame-independent consequences result. The Ehrenfest paradox therefore shows that the Lorentz contraction, although absent in each local co-moving description, has measurable global consequences when the local descriptions cannot be made mutually compatible.


%Each individual Lorentz contraction can be made to disappear by passing to the instantaneous rest frame of the corresponding material element. In the case of uniform translation, these local rest frames fit together into a single global co-moving inertial frame, so the contraction can be eliminated for the body as a whole. In the case of a rotating disk, however, the instantaneous rest frames of the different elements cannot be assembled into one global co-moving inertial frame. The contractions may therefore be removed locally, but not all at once. The failure of this local elimination to globalize is what gives rise to the measurable stresses or strains in the disk.
%
%Even sharper for the response letter:
%The point is not that Lorentz contraction is real merely because it cannot be eliminated in all frames. Rather, each local contraction can be eliminated by adopting the co-moving frame of the relevant element. What the Ehrenfest disk shows is that, for an extended rotating body, these local co-moving descriptions cannot be made mutually compatible as one global rest description. The effect therefore disappears locally but not globally, and this failure has measurable consequences.


%
%In fact, one could say that the twin paradox is the temporal counterpart of the Ehrenfest paradox: in one case no global co-moving inertial space exists for the rotating body; in the other no global co-moving inertial time frame exists for the accelerated twin. The invariant observable (differential aging or stress) emerges precisely from that failure of globalization.

%The case of the rotating disk is used by Einstein in order to show that in this case the effect cannot be even be removed by the co-moving observer
%
%In uniform translation, the description in the co-moving frame and the description in a non-co-moving inertial frame remain globally equivalent: all collinear parts of the rod acquire the same Lorentz factor, so that the ratio between the whole and its parts is preserved. In rotation this equivalence breaks down. Each material element has an instantaneous rest frame, but these local rest frames cannot be assembled into a single global co-moving inertial frame. The contraction of rim elements is tangential, whereas radial elements are orthogonal to the motion and are not contracted in the same way. Thus, if the disk were required to preserve its original proper circumference and radius during spin-up, the inertial-frame circumference would have to be \(C'/\gamma\), while the radius would remain \(R'\), contradicting the Euclidean relation \(C=2\pi R\). The Ehrenfest paradox therefore shows that the Lorentz contraction, although absent in each local co-moving description, has measurable global consequences when the local descriptions cannot be made mutually compatible.
\end{reply}

\begin{point}\label{vanish}
The transverse Doppler effect concerns the claim is that the frequency of emission of a moving ion, the emission being in the transverse direction, is different from the proper frequency as defined relative to the co-moving frame (which is the same as the proper frequency defined relative to K). It is arguably misleading to say that the effect "vanishes" relative to the co-moving frame. The effect might be described as a relation between two inertial frames, as is brought out on p. 16. Better put, it is a relation between proper frequency relative to K and that related to an ion moving relative to K. In the author's scenario, relative to K' this relation does not exist— there is a single "clock" and one cannot even define the Doppler shift, and so it cannot have zero value (if that is what the author means by "vanishing"). However, in the reversed scenario in which the observer at rest in K' is looking at emitting ions at rest relative to K, then the transverse Doppler shift
is defined and the effect is the same as in the original case.
\end{point}


\begin{reply}

%I use “disappears” or “vanishes” in the ordinary operational sense familiar from Doppler spectroscopy. A Doppler effect is absent when the spectrograph records no displacement of the spectral line from its rest position. Thus the classical longitudinal Doppler effect disappears either when there is no relative motion along the line of sight or when the observation is arranged transversely. Analogously, the transverse Doppler shift disappears for a spectrograph co-moving with the emitting ion, since it records the proper frequency \(\nu_0\). This does not mean that the transverse Doppler effect is an intrinsic property of the ion which takes the value zero in its rest frame; it means that the spectral displacement is not observed in the co-moving experimental arrangement.

I did not realize that my use of the term \s{vanish} was indeed ambiguous. By saying that the effect \s{vanishes} for the co-moving observer, I use the term in an \s{operational} sense. The relevant question is what is recorded by measuring devices at rest with the system under investigation. A rod measured by co-moving rods has its proper length; a clock compared with co-moving clocks has its proper rate; an emitting ion observed by a co-moving spectrograph displays its proper frequency. In such an arrangement, the observable difference that defines the effect in a non-co-moving frame is not registered.

Thus, in the transverse Doppler case, the claim is not that the Doppler shift is an intrinsic property of the ion which somehow takes the value zero in its rest frame. Rather, a spectrograph co-moving with the emitting ion records the normal rest line, \(\nu=\nu_0\), whereas a spectrograph in transverse relative motion records the shifted line, \(\nu=\gamma^{-1}\nu_0\). This is the precise sense in which the shift \s{vanishes} for the co-moving observer. It is absent in the co-moving observational arrangement, but experimentally measurable in the non-co-moving one.



\end{reply}


\begin{point}
The second point I want to make is more important. The author attempts to show, through an analysis of an important historical episode related to the testing of time dilation, that there is a confusion in the claim by Brown and others that such a "kinematical" effect (to use Einstein's original wording) has a dynamical underpinning, despite Einstein's belated recognition, endorsed by the author, that rods and clocks are physical systems governed by dynamical laws. Essentially, the author's argument is this. A dynamical account of, specifically, the transverse Doppler effect in special relativity, under plausible assumptions, would explain why all atomic clocks are identical when compared in the same rest frame, but no more. Explaining the transverse frequency shift requires, according to the author, more than that: it further requires time dilation regarded as a kinematical, not dynamical effect (pp. 3, 24).

Here are some objections to this view. 

Bell showed in his famous 1976 paper on teaching special relativity that a purely dynamical (pre-quantum) model of a moving single-electron atom has time dilation unequivocally emerging from dynamical assumptions (largely Maxwell electrodynamics) in this specific model. Recall his point: the special merit of this approach "is to drive home the lesson that the laws of physics in any one reference frame account for all physical phenomena, including the observations of moving observers" (Bell 1976). The author's appeal to irreducible kinematics in fully accounting for the transverse Doppler effect seems to be in direct conflict with this important insight.

Furthermore, Bell makes a nice analogy with the parallel roles of thermodynamics (which gave Einstein the methodological template for his 1905 relativity paper) and statistical mechanics. He wrote: 

"If you are, for example, quite convinced of the second law of thermodynamics, of the increase of entropy, there are many things that you can get directly from the second law which are very difficult to get directly from a detailed study of the kinetic theory of gases, but you have no excuse for not looking at the kinetic theory of gases to see how the increase of entropy actually comes about. In the same way, although Einstein's theory of special relativity would lead you to expect the FitzGerald contraction, you are not excused from seeing how the detailed dynamics of the system also leads to the FitzGerald contraction." (Bell 1976, as referenced in the paper.)

Of course the same can be said of time dilation. The key point is this. Naturally, Einstein used time dilation as established in the "kinematical" part of his 1905 paper in his treatment of the transverse Doppler effect. This was because the full dynamical theory of moving Ions and their emission properties was not available to him, and even today would be very complicated. The beauty of the kinematical approach was that it did not rely on the fine-grained details of the dynamics, as the author acknowledges. It is analogous to the way thermodynamics leads to an entropy non-decrease law without specifying all the complicated details that would be required in a purely mechanical analysis of the relevant system. But as the current author concedes, special relativity is a set of constraints on theories. Indeed, it is much the way that thermodynamics is. Because something is hard to prove in what Einstein called a constructive (dynamical) theory of matter and radiation — even if
it available —, while at the same amenable to a simple principle-theory treatment, does not mean that the latter is fundamental, or in principle irreplaceable. The author takes the lesson of the transverse Doppler effect to be that "kinematic and dynamics, from a strictly logical point of view, cannot be confronted with experience separately but only as a whole. From the point of view of explanation, however, the division of labour within this whole remains intact." (pp. 23, 24).  The analogous, erroneous, claim would be that certain phenomena in the behaviour of gases, for example, can only be explained in principle by an irreducible combination of thermodynamics and statistical mechanics.
\end{point}

\begin{reply}
O thank the referee for raising this important point. The point of the paper is not to deny that Bell’s program provides an independent theoretical motivation for a dynamical approach to relativity. Rather, it is to question a specific historical inference: namely, that Einstein’s appeal to a dynamical theory of rods and clocks shows that he regarded time dilation itself as standing in need of a dynamical explanation. On the interpretation defended in the paper, Einstein’s point was different. A dynamical theory of matter is needed to account for the physical reliability of clocks, and in the case discussed here for the spectral identity of identical atoms. Once this identity has been secured, however, the transverse Doppler shift is explained kinematically, as a consequence of the Lorentz transformation of time. In this sense, the effect is \s{apparent} in Einstein’s terminology, though empirically real. A dynamical explanation of the shift itself would seem to require the Lorentz--Ives strategy of treating the moving atom as undergoing a real change of intrinsic frequency. In the new version of the paper, I point out explicitly that Bell’s approach may offer a third possibility, but it is an independent philosophical development, not evidence for the historical claim that Einstein himself sought a dynamical explanation of time dilation.

%The paper does not deny that one can construct a detailed model calculation in which a material system, described from a single inertial frame, exhibits the Lorentz-contracted or time-dilated behavior expected by special relativity. Bell’s example is valuable precisely because it shows how relativistic behavior can be realized by concrete physical systems governed by appropriate dynamical laws.
%
%The point at issue, however, is how such a calculation is to be interpreted. 
%
%\begin{itemize}
%\item a If Bell’s procedure is read as a genuinely constructive explanation of time dilation in the strong dynamical sense, then the motion of the atom relative to some privileged state of rest must make the moving atom dynamically different from an atom at rest. This is close to the neo-Lorentzian strategy pursued by Herbert E. Ives in connection with the Ives-Stilwell experiment on the rate of a moving atomic clock. On that interpretation, the transverse Doppler effect is explained by a real dynamical alteration of the moving atomic system. But this is not the constructive program usually intended by defenders of a dynamical approach within special relativity, since it reintroduces precisely the preferred-frame structure that Einstein’s theory was designed to avoid.
%
%\item If, by contrast, Bell’s model is interpreted within Einsteinian special relativity, then the relevant dynamical laws are Lorentz covariant. A moving atom is not dynamically different from an atom at rest; it is the Lorentz-transformed counterpart of the same kind of physical solution. In that case the dynamics explains why stable atomic clocks exist, why atoms of the same kind have the same proper frequency when compared at rest, and why such systems can serve as clocks. But the transverse Doppler factor itself follows from the Lorentz-transformational relation between the proper time of the emitting ion and the time coordinate of the non-co-moving observational frame. The dynamical model realizes the relativistic kinematics; it does not replace it with an independent dynamical alteration of the clock.
%\end{itemize}
%
%This is also how we understand the analogy with thermodynamics and statistical mechanics. The point is not that principle-level explanations are in principle irreplaceable by constructive accounts. Rather, the point is that a constructive account succeeds only by showing how the concrete system instantiates the relevant general constraints. In the present case, any acceptable special-relativistic dynamics of matter and radiation must be Lorentz covariant. Once that condition is imposed, the transverse Doppler effect follows generally for any suitable clock with a stable proper frequency. The detailed dynamics may explain the existence and stability of the clock, but the occurrence of the universal \(\gamma\)-factor is fixed by the Lorentz-covariant structure itself.
%
%We have therefore revised the manuscript to avoid suggesting that a Bell-style dynamical calculation is impossible or unimportant. The intended claim is narrower: within Einsteinian special relativity, a dynamical account of the atom explains the identity and stability of the atomic clock, while the transverse Doppler shift is explained by the Lorentz transformation relating the clock’s proper time to the time coordinate of a non-co-moving frame. A stronger dynamical explanation of the shift itself is available only by moving toward a Lorentzian or Ives-style interpretation, according to which moving atoms undergo a real dynamical modification relative to a preferred frame. I do not see the third alternative of a dynamical explanation within relativity, without retuirn to the ether theory.
\end{reply}

\begin{point}
It stated that the Lorentz transformations make their ``perfect symmetry in $x$ and $t$ explicit''. Is it perfect?
\end{point}

\begin{reply}
I agree that the formulation was too strong. I have replaced it with a more cautious statement: the Lorentz transformations make explicit the reciprocal mixing of spatial and temporal coordinates, especially when the time coordinate is written as $ct$.
\end{reply}

\begin{point}
It is claimed that Einstein was the first to see time dilation as a genuinely new physical effect. But Larmor and Lorentz independently had a clear premonition of time dilation in 1897 and 1899 respectively, though it is not clear that they thought of it as a universal property of ideal clocks.
\end{point}

\begin{reply}
I thank the referee for this correction. I have qualified the claim. The revised text now acknowledges that Larmor and Lorentz had obtained related formal results before Einstein. The point I intended to make is narrower: Einstein’s distinctive contribution was to interpret clock retardation as a universal consequence of the new kinematics, applying to any suitable clock, rather than as a feature of a particular electromagnetic model.
\end{reply}

\begin{point}
The coordinate time interval in $K$ between two successive ticks of a moving clock cannot be measured by a single clock at rest in $K$, since the ticks occur at different spatial points in $K$. This risks inconsistency with the claim that the moving clock runs ``more slowly'' than an identical clock at rest in $K$.
\end{point}

\begin{reply}
I agree. I have revised the relevant passage to make explicit that $\Delta t$ is read off by the synchronized lattice of clocks at rest in $K$, whereas $\Delta t'$ is the proper time measured by the single clock at rest in $K'$. I have also qualified the potentially misleading phrase that the moving clock ``runs more slowly'' with a formulation in terms of the longer coordinate-time interval assigned in $K$ to successive ticks of the moving clock.
\end{reply}

\begin{point}
The manuscript seems to treat the invariance of the number $N$ of oscillations as a substantive assumption based on the relativity principle. But $N$ is a count and is independent of all coordinate frames by its nature.
\end{point}

\begin{reply}
I agree that the previous wording was misleading. I have revised the manuscript throughout to distinguish the trivial invariance of the count $N$ from the substantive physical premise. The point is not that $N$ itself is made invariant by the relativity principle, but that uniform motion is assumed not to alter the intrinsic oscillatory process of the emitter. Thus identical ions, when compared at rest, possess the same proper frequency $\nu_0=N/\Delta t'$. The transverse Doppler effect then arises because the same count $N$ is referred to different time intervals, $\Delta t'$ and $\Delta t=\gamma\Delta t'$, in the two inertial descriptions.
\end{reply}

\begin{point}
Ritz is described as criticizing the abandoning of absolute time and rigid bodies in the Lorentz--Einstein theory. But this theory does not abandon idealized rigid bodies.
\end{point}

\begin{reply}
I agree. I have revised the wording to clarify that Ritz was criticizing the abandonment of absolute time and the classical notion of absolutely rigid bodies, not the use of idealized rods or rigid bodies as measuring instruments within the relativistic framework. More specifically, Ritz had in mind something similar to the contrast between Abraham's rigid electron and the deformable electron of Lorentz's and Einstein's theories.
\end{reply}

\begin{point}
Typos: first line p.~16; third line final paragraph p.~18; equation in last line of first paragraph of section~4, p.~20.
\end{point}

\begin{reply}
I thank the referee for pointing these out. I have corrected these typographical errors in the revised manuscript.\todo{??}
\end{reply}


\reviewersection

\begin{point}
The present submission contributes to the literature on the relation between
dynamical and kinematical considerations in special relativity.  Specifically,
it disputes a narrative to the effect that Einstein initially treated time
dilation as a merely apparent coordinate artifact but sought a dynamical
explanation on which time dilation reflects real changes, thereby suggesting a
precursor to today's dynamical view.  To challenge this narrative, the paper
reconstructs in detail the role of the kinematical/dynamical distinction in
Einstein's treatment of the transverse Doppler effect.  It argues that Einstein
was interested not in a dynamical explanation of the transverse Doppler effect
but rather in a dynamical backing for assumptions going into experimental
testing of the effect.  

This paper is timely because it concerns debates over kinematical and dynamical
explanation in special relativity.  The topic is also relevant to those with
historical interests in the special theory of relativity.  The historical
analysis exhibits detailed engagement with the primary sources and substantial
work in synthesizing them.  It successfully makes its conceptual case on the
basis of this analysis.

Minor typos:
2: "as [shown] by the stresses in Ehrenfest's disk"
9: "[in] November of 1905"
23: "which are in turn required to invariant with [respect to] the Lorentz
transformations"
\end{point}

\begin{reply}
I have corrected the remaining typos where this phrase still appeared in the revised version.
\end{reply}






\end{document}